\section{Architecture and Foundations}\label{sec:architecture}

NEOS v1.0 is a \textbf{specification}---37 markdown files that define
an interactive runtime environment running \emph{on} an LLM.  The LLM
is the virtual machine.  NEOS is the operating system.  Neural fields
are the computational substrate.

This section presents the governing dynamics and the three
foundational pillars on which the entire system rests.

%% ════════════════════════════════════════════════════════════════════
\subsection{The Master Equation}\label{sec:master-equation}

All dynamics in NEOS are governed by a single partial differential
equation:

\begin{eqbox}
\begin{equation}\label{eq:master-pde}
\frac{\partial A}{\partial t}
  = \underbrace{-\lambda\, A(x)}_{\text{Decay}}
  + \underbrace{\alpha \int K(x,y)\, A(y)\, dy}_{\text{Resonance}}
  + \underbrace{\iota(x,t)}_{\text{Injection}}
\end{equation}
\end{eqbox}

\noindent
This equation belongs to the family of neural field
equations introduced by Amari~\cite{amari1977dynamics} and
Wilson--Cowan~\cite{wilson1972excitatory}, adapted here for a semantic
rather than neural substrate.  Each term plays a distinct cognitive
role:

\begin{description}[leftmargin=2.2cm, labelwidth=2cm, style=nextline]

\item[\boldmath$-\lambda\, A(x)$\unboldmath\quad Decay.]
  Unreinforced ideas fade exponentially.  This is not a defect but a
  design principle: \emph{forgetting is a feature, not a bug}.  In the
  absence of resonance support, every pattern loses activation at rate
  $\lambda$, ensuring that the field does not accumulate noise
  indefinitely.  After the decay phase of each cycle, the activation
  update is $A \leftarrow A \times (1 - \lambda)$.

\item[\boldmath$\alpha \int K(x,y)\, A(y)\, dy$\unboldmath\quad Resonance.]
  Related ideas strengthen each other.  The resonance kernel
  $K(x,y)$ quantifies the semantic, logical, and contextual affinity
  between patterns located at positions $x$ and $y$ in the field.
  When two patterns resonate, each contributes to the other's
  activation proportionally to their affinity and their current
  strength.  The amplification factor $\alpha$ controls how aggressively
  mutual reinforcement operates.  This integral is the ``hearing''
  mechanism: a pattern survives not because it is loud, but because it
  is \emph{heard} by others.

\item[\boldmath$\iota(x,t)$\unboldmath\quad Injection.]
  The user adds new ideas at any time via the \texttt{nf inject}
  command.  Each injection introduces a fresh signal---a pattern with
  an initial activation strength---into the field at position $x$ and
  time $t$.  The injected pattern immediately begins interacting with
  all existing patterns through the resonance integral.

\end{description}

\noindent
From this single equation, \emph{all} of NEOS emerges: attractor
formation, coherence dynamics, collapse, versioning, and multi-field
orchestration.

\paragraph{Parameters.}
Four scalar parameters define the \textbf{cognitive style} of a field.
A security analysis benefits from low $\lambda$ (do not forget
threats), high $\alpha$ (cluster related threats fast), and high $\tau$
(only credible concerns survive).  A creative brainstorm benefits from
high $\lambda$ (rapid turnover), moderate $\alpha$, and low $\tau$
(let weak ideas live).  Same engine; different cognitive personality;
tuned with four numbers.

\begin{table}[htbp]
\centering
\caption{Master-equation parameters and their effects on field
  dynamics.}
\label{tab:parameters}
\small
\begin{tabular}{@{}cllcl@{}}
\toprule
\thead{Symbol} & \thead{Name} & \thead{Default}
  & \thead{Range} & \thead{Effect} \\
\midrule
$\lambda$ & Decay rate     & 0.05 & $[0,\,1]$
  & Higher $=$ faster forgetting \\
$\alpha$  & Amplification  & 0.30 & $[0,\,1]$
  & Higher $=$ stronger resonance \\
$\tau$    & Threshold      & 0.40 & $[0,\,1]$
  & Below this, patterns marked dormant \\
$\sigma$  & Bandwidth      & 0.50 & $[0,\,\infty)$
  & Semantic reach of each pattern \\
\bottomrule
\end{tabular}
\end{table}

%% ════════════════════════════════════════════════════════════════════
\subsection{Neural Fields --- The Substrate}\label{sec:pillar-fields}

Classical prompt engineering represents information as a linear
sequence of tokens with a hard context-window boundary.  NEOS replaces
this with \textbf{continuous activation landscapes}---neural fields in
which information is encoded not as discrete tokens but as spatially
extended activation patterns~\cite{amari1977dynamics}.

Five principles govern the substrate:

\begin{enumerate}[label=\textbf{\arabic*.}, leftmargin=2em]
\item \textbf{Continuity.}  Activation is a continuous function over a
  semantic space, not a discrete list.  Nearby meanings have nearby
  activations.

\item \textbf{Resonance.}  Patterns influence one another through the
  kernel $K(x,y)$.  If two ideas are semantically proximate, injecting
  one raises the activation of the other.

\item \textbf{Persistence.}  Information persists through resonance
  support, not through explicit inclusion in a prompt.  A pattern
  endures as long as the field sustains it.

\item \textbf{Entropy.}  The field tracks its own information-theoretic
  entropy, $H = -\sum p_i \log_2 p_i$, providing a scalar measure of
  how concentrated or diffuse the current state is.

\item \textbf{Boundary Dynamics.}  Patterns near the activation
  threshold $\tau$ are in a critical zone: small resonance boosts can
  save them, small decays can extinguish them.  This boundary behaviour
  produces the field's sensitivity to context.
\end{enumerate}

\paragraph{Why fields beat prompts.}
In a prompt, information persists only if explicitly included.  Context
window overflow silently destroys old information.  In a neural field,
information persists through \emph{resonance}---if a pattern connects
to others that keep it alive, it endures.  Relationships emerge
naturally from semantic proximity.  New input interacts with the
\emph{entire} field, not just recent tokens.

%% ════════════════════════════════════════════════════════════════════
\subsection{Symbolic Reasoning --- The Structure}\label{sec:pillar-symbolic}

Without structure, fields are beautiful but unusable---like having a
piano with no keys.  The \texttt{nf} command set provides the
\textbf{opcodes of cognitive computing}: a symbolic language for field
manipulation that is both human-readable and algebraically precise.

Core operations include:
\texttt{inject} (add a pattern),
\texttt{amplify} (increase activation),
\texttt{attenuate} (decrease activation),
\texttt{cycle} (advance dynamics),
\texttt{collapse} (resolve superposition), and
\texttt{commit} (snapshot state).
Each command has a precise mathematical meaning in terms of the master
equation~\eqref{eq:master-pde}.

Two reference conventions provide the symbolic glue:

\begin{itemize}[leftmargin=1.5em]
\item \textbf{Pattern references} (\texttt{@name}) select an
  individual pattern in the field.
\item \textbf{Field references} (\texttt{\$field}) select an entire
  named field in multi-field configurations.
\end{itemize}

\noindent
Operations compose algebraically.  For example, applying
\texttt{amplify(@a, 1.5)} followed by \texttt{attenuate(@a, 0.5)}
returns the pattern to a well-defined intermediate activation
that can be computed from the field state alone.

%% ════════════════════════════════════════════════════════════════════
\subsection{Quantum Semantics --- The Interpretation}%
\label{sec:pillar-quantum}

Before collapse, meaning exists in
\textbf{superposition}~\cite{busemeyer2012quantum}.  The field holds
multiple possible interpretations simultaneously, each weighted by
activation strength.  This is not a loose metaphor: the mathematical
structure of NEOS's collapse operation mirrors key features of quantum
measurement~\cite{pothos2013can}.  We adopt the term ``quantum
semantics'' from the quantum cognition literature; the formalism draws
on quantum probability theory as a model of cognitive phenomena, not on
quantum physics.

\paragraph{Non-commutativity.}
Injecting pattern $A$ and then pattern $B$ is \emph{not} the same as
injecting $B$ and then $A$.  Each injection alters the field that
receives the next one.  Formally, if $I_A$ and $I_B$ denote the
injection operators, then in general $I_A \circ I_B \neq I_B \circ
I_A$ because the resonance integral couples each new pattern to the
full existing field state.  Order of inquiry shapes the space of
possible conclusions.

\paragraph{Observer-dependent collapse.}
When the operator issues \texttt{nf collapse}, the superposition
resolves into a definite output---but the result depends on the
chosen strategy:

\begin{itemize}[leftmargin=1.5em]
\item \textbf{Attractor}: select the dominant attractor basin.
\item \textbf{Threshold}: retain all patterns above a specified
  activation.
\item \textbf{Weighted}: blend all patterns proportionally to their
  activation.
\item \textbf{Sample}: probabilistically sample from the activation
  distribution.
\end{itemize}

\noindent
The observer matters.  The strategy you choose shapes the output you
get.  The same field can collapse differently under different
strategies, just as the same quantum state yields different
measurement outcomes under different observables.

%% ════════════════════════════════════════════════════════════════════
\subsection{The Three Pillars in Concert}\label{sec:pillars-unified}

\begin{figure}[H]
\centering
\begin{tikzpicture}[
    >=Stealth
  ]

  %% --- Three pillars ---
  \node[draw=neosteal!60, fill=neosteal!6, minimum width=4.3cm, minimum height=2.4cm, align=center, font=\small, rounded corners=6pt] (nf) at (-5,0)
    {\textbf{Neural Fields}\\[4pt]
     {\footnotesize THE SUBSTRATE}\\[2pt]
     {\scriptsize Continuous activation}\\
     {\scriptsize landscapes}};

  \node[draw=neosviolet!60, fill=neosviolet!6, minimum width=4.3cm, minimum height=2.4cm, align=center, font=\small, rounded corners=6pt] (sr) at (0,0)
    {\textbf{Symbolic Reasoning}\\[4pt]
     {\footnotesize THE STRUCTURE}\\[2pt]
     {\scriptsize Composable command}\\
     {\scriptsize algebra (\texttt{nf})}};

  \node[draw=neosorange!60, fill=neosorange!6, minimum width=4.3cm, minimum height=2.4cm, align=center, font=\small, rounded corners=6pt] (qs) at (5,0)
    {\textbf{Quantum Semantics}\\[4pt]
     {\footnotesize THE INTERPRETATION}\\[2pt]
     {\scriptsize Superposition \&}\\
     {\scriptsize observer collapse}};

  %% --- Unification bar ---
  \node[rectangle, rounded corners=3pt, draw=neosblue!60, fill=neosblue!12,
        minimum width=14.8cm, minimum height=0.7cm,
        font=\small\bfseries, text=neosblue!70!black]
    (bar) at (0,-2.2) {NEOS: The Unification};

  %% --- Connections ---
  \draw[thick, neosblue!40] (nf.south)  -- (bar.north -| nf.south);
  \draw[thick, neosblue!40] (sr.south)  -- (bar.north);
  \draw[thick, neosblue!40] (qs.south)  -- (bar.north -| qs.south);
\end{tikzpicture}
\caption{The three pillars of NEOS.  Neural fields provide the
  computational substrate; symbolic reasoning provides the structural
  language; quantum semantics provides the interpretive framework.
  NEOS unifies all three into a single coherent system.}
\label{fig:three-pillars}
\end{figure}

Each pillar is necessary; together they make cognitive computing
possible.  Remove neural fields and the system cannot \emph{represent}
meaning.  Remove symbolic reasoning and it cannot \emph{manipulate}
meaning.  Remove quantum semantics and it cannot \emph{interpret}
meaning.  Figure~\ref{fig:three-pillars} shows the relationship.

\paragraph{Collapse visualised.}
Figure~\ref{fig:quantum-collapse} illustrates the collapse operation.
Before collapse, the field contains multiple patterns in superposition,
each with a distinct activation level.  After the operator issues
\texttt{nf collapse}, the superposition resolves into structured
output---the dominant interpretation crystallised from the full field
state.

\begin{figure}[H]
\centering
\begin{tikzpicture}[>=Stealth]

  %% ===== BEFORE (left panel) =====
  \node[font=\small\bfseries, text=gray!60!black] at (-3.8,2.6)
    {Before collapse};

  \draw[gray!30, rounded corners=6pt]
    (-6.2,-1.8) rectangle (-1.4,2.2);

  % Multiple patterns (circles of varying size/shade)
  \fill[neosblue!25]   (-5.0, 1.2) circle (0.45);
  \fill[neosblue!50]   (-3.6, 1.5) circle (0.55);
  \fill[neosblue!15]   (-2.4, 0.8) circle (0.30);
  \fill[neosblue!40]   (-4.8,-0.2) circle (0.50);
  \fill[neosblue!60]   (-3.2, 0.0) circle (0.65);
  \fill[neosblue!10]   (-2.2,-0.5) circle (0.25);
  \fill[neosblue!35]   (-4.0,-1.0) circle (0.40);
  \fill[neosblue!20]   (-2.6,-1.2) circle (0.35);

  % Labels
  \node[font=\tiny, text=gray] at (-5.0, 1.2) {@a};
  \node[font=\tiny, text=white] at (-3.6, 1.5) {@b};
  \node[font=\tiny, text=gray] at (-2.4, 0.8) {@c};
  \node[font=\tiny, text=gray] at (-4.8,-0.2) {@d};
  \node[font=\tiny, text=white] at (-3.2, 0.0) {@e};
  \node[font=\tiny, text=gray] at (-4.0,-1.0) {@f};

  %% ===== ARROW =====
  \draw[->, ultra thick, neosblue!50] (-0.9,0.2) -- (0.9,0.2)
    node[midway, above, font=\small\bfseries, text=neosblue!60!black]
      {\texttt{nf collapse}};

  %% ===== AFTER (right panel) =====
  \node[font=\small\bfseries, text=gray!60!black] at (3.8,2.6)
    {After collapse};

  \draw[gray!30, rounded corners=6pt]
    (1.4,-1.8) rectangle (6.2,2.2);

  % Single dominant circle
  \fill[neosblue!65] (3.8,0.2) circle (1.2);
  \node[font=\small\bfseries, text=white, align=center] at (3.8,0.2)
    {Structured\\Output};

  % Faint remnants
  \fill[neosblue!8] (5.4, 1.4) circle (0.15);
  \fill[neosblue!8] (5.6, 0.0) circle (0.12);
  \fill[neosblue!8] (5.2,-1.0) circle (0.10);

\end{tikzpicture}
\caption{Field collapse (quantum-inspired semantics).  \textbf{Left:} the field in
  superposition---multiple patterns with varying activation levels.
  \textbf{Right:} after \texttt{nf collapse}, the dominant
  interpretation crystallises into structured output.  Residual
  patterns fade to near-zero activation.}
\label{fig:quantum-collapse}
\end{figure}
